\documentclass[10pt,a4paper]{amsart}
\usepackage[margin=19mm]{geometry}
\usepackage{amsmath,amssymb,mathtools}
\usepackage[scr=rsfs,cal=euler]{mathalfa}
\usepackage{xfrac}
\usepackage[unicode,hidelinks,bookmarks=false]{hyperref}
\newcommand{\ie}{i.e.\ }
\newcommand{\cf}{cf.\ }
\newcommand{\resp}{respectively\ }
\newcommand{\Subsection}{\subsection}
\providecommand{\nobreakdash}{\nobreak\hbox{-}}
\setlength{\emergencystretch}{2em}
\multlinegap0pt
\usepackage{amssymb}
\usepackage[shortlabels]{enumitem}
\newtheorem{proposition}{Proposition}
\newcommand{\Spec}{\mathrm{Spec}}
\newcommand{\bbT}{\mathbb{T}}
\title{Spectral--Geometric Deformations of Function Algebras on Manifolds}
\author{Amandip Sangha}
\date{Source: arXiv:2510.24184v3 (CC BY 4.0)}
\begin{document}
\maketitle

\noindent\emph{Source and licence.} Spectral--Geometric Deformations of Function Algebras on Manifolds. Version arXiv:2510.24184v3 (CC BY 4.0), distributed by arXiv under the Creative Commons Attribution 4.0 International licence, http://creativecommons.org/licenses/by/4.0/, as recorded in the arXiv licence metadata for this exact version and shown on the abstract page https://arxiv.org/abs/2510.24184v3. Full licence notice: \texttt{licenses/arxiv-2510.24184-CC-BY-4.0.txt}.\par\smallskip

\noindent\textit{Editorial scope.} Counted unit: the article's proposition prop:separable\_twists\_are\_coboundaries (source0.tex lines 759--771). The article's complete original proof of it is delivered with it (source0.tex lines 773--821, one \texttt{proof} environment of 49 lines). No mathematics is written, repaired or re-derived: the statement and the proof are the article's own lines, in the article's own notation, and no retained line is substituted or re-flowed.  Retained uncounted as the source-local context the proof needs: lines 702--717.  Nothing was omitted from any retained span, so the package carries no omission record.  No retained line contains a citation, so the wrapper carries no bibliography and no bibliography entry.  No retained line contains a figure environment, an image-inclusion command or drawing code, so no image is delivered and the package declares no asset dependency.  Editorial formatting only, in addition: an AMS \texttt{amsart} preamble that restates the article's own declarations and macros used by the retained lines, namely the environments \texttt{proposition} on separate counters, together with the standard packages those lines need.  The retained environments print in this document as Proposition 1 in order of appearance, whereas the article prints the same items with the numbering of its own full document.  The complete native source of the article as distributed in the arXiv e-print for this exact version is bundled unchanged as \texttt{sources/187934/source0.tex}.
\begin{proposition}[Gauge equivalence by a spectral coboundary]\label{prop:gauge_equivalence}
Let $\omega=(\omega^\nu_{\lambda\mu})$ and $\omega'=(\omega'{}^\nu_{\lambda\mu})$ be two scalar twisting systems.
Assume there exists a function $a:\Spec(\Delta)\to\bbT$ such that for all $\lambda,\mu,\nu$,
\begin{equation}\label{eq:coboundary_relation}
\omega'{}^\nu_{\lambda\mu} \;=\; a(\lambda)a(\mu)\overline{a(\nu)}\,\omega^\nu_{\lambda\mu}.
\end{equation}
Define the diagonal unitary $A:=\sum_\lambda a(\lambda)P_\lambda$ on $L^2(M)$.

Then for all $f,g\in E_{\mathrm{fin}}$,
\[
A(f\star_{\omega'} g) \;=\; (Af)\star_\omega (Ag).
\]
Thus $A$ is a linear bijection $E_{\mathrm{fin}}\to E_{\mathrm{fin}}$ intertwining the core products (as maps into $L^2(M)$).
Moreover, if $\star_\omega$ and $\star_{\omega'}$ extend continuously to $H^s(M)$ for some $s$,
then $A$ extends to a bounded algebra isomorphism on $H^s(M)$ intertwining the extended products.
\end{proposition}

\begin{proposition}[Rigidity of separable twists]\label{prop:separable_twists_are_coboundaries}
Assume $\omega=(\omega^\nu_{\lambda\mu})$ has the separable form
\[
\omega^\nu_{\lambda\mu}=p(\lambda)\,q(\mu)\,\overline{r(\nu)},
\qquad p,q,r:\Spec(\Delta)\to\bbT.
\]
Assume moreover that $\star_\omega$ is unital with unit $1\in E_0$, i.e.\ $f\star_\omega 1=f=1\star_\omega f$ for all $f\in E_{\mathrm{fin}}$.
Then $\omega$ is a spectral coboundary: there exists $a:\Spec(\Delta)\to\bbT$ with $a(0)=1$ such that
\[
\omega^\nu_{\lambda\mu}=a(\lambda)a(\mu)\overline{a(\nu)}.
\]
Consequently, $\star_\omega$ is gauge-trivial in the sense of Proposition~\ref{prop:gauge_equivalence}.
\end{proposition}

\begin{proof}
Unitality implies, for each $\lambda$, that $f_\lambda\star_\omega 1=f_\lambda$ for all $f_\lambda\in E_\lambda$.
Since $1\in E_0$ and $f_\lambda\cdot 1=f_\lambda\in E_\lambda$, we have
\[
P_\nu(f_\lambda\cdot 1)=\delta_{\nu\lambda}\,f_\lambda,
\]
so in the defining formula for $P_\nu(f_\lambda\star_\omega 1)$ only the term with $\nu=\lambda$ survives, and hence $\omega^\lambda_{\lambda 0}=1$.
By the separable form this reads
\[
p(\lambda)\,q(0)\,\overline{r(\lambda)}=1\qquad\text{for all }\lambda,
\]
hence $p(\lambda)=\overline{q(0)}\,r(\lambda)$.

Similarly, $1\star_\omega g_\mu=g_\mu$ for $g_\mu\in E_\mu$ gives $\omega^\mu_{0\mu}=1$, i.e.
\[
p(0)\,q(\mu)\,\overline{r(\mu)}=1\qquad\text{for all }\mu,
\]
so $q(\mu)=\overline{p(0)}\,r(\mu)$.

Therefore
\[
\omega^\nu_{\lambda\mu}
=
p(\lambda)\,q(\mu)\,\overline{r(\nu)}
=
\overline{q(0)p(0)}\ r(\lambda)r(\mu)\overline{r(\nu)}.
\]
Now $1\star_\omega 1=1$ implies $\omega^0_{00}=1$, so
\[
1=\omega^0_{00}=\overline{q(0)p(0)}\ r(0),
\qquad\text{hence }\overline{q(0)p(0)}=\overline{r(0)}.
\]
Thus
\[
\omega^\nu_{\lambda\mu}=\overline{r(0)}\ r(\lambda)r(\mu)\overline{r(\nu)}.
\]
Define $a:\Spec(\Delta)\to\bbT$ by $a(\lambda):=\overline{r(0)}\,r(\lambda)$.
Then $a(0)=\overline{r(0)}r(0)=1$ and
\[
a(\lambda)a(\mu)\overline{a(\nu)}
=
(\overline{r(0)}r(\lambda))(\overline{r(0)}r(\mu))\overline{\overline{r(0)}r(\nu)}
=
\overline{r(0)}\ r(\lambda)r(\mu)\overline{r(\nu)}
=
\omega^\nu_{\lambda\mu},
\]
as required.
\end{proof}

\end{document}
